% ----------------------------------------------------------------------
%
%                         Computer Science VII
%
%                   http://ls7-www.cs.uni-dortmund.de
%
%   questions and representations: info@ls7.cs.uni-dortmund.de
%
%   status: 20.12.2017
%
% ----------------------------------------------------------------------

\RequirePackage{ifthen}
\newcommand \Thesistyp{Masterarbeit}
\newcommand \Author{Jessica B\"uhler}
\newcommand \Thesistitle{Zeit-effizientes Training von Convolutional Neural Networks}
\newcommand \FirstSupervisor{Prof.~Dr.~Heinrich~M\"uller}
\newcommand \SecondSupervisor{M.Sc.~Matthias~Fey}
\newcommand \FirstChair{Lehrstuhl VII}
\newcommand \FirstChairTitle{Informatik}


% -----------------------------------------------------------------------------------------
% Option: Zweiter Lehrstuhl
\newboolean{boolNoSecondChair}
\setboolean{boolNoSecondChair}{true} % boolNoSecondChair==false: second chair involved
\ifthenelse{\boolean{boolNoSecondChair}}{
\newcommand \SecondChair{}
\newcommand \SecondChairTitle{}
}{
\newcommand \SecondChair{Computer Science XII}
\newcommand \SecondChairTitle{Embedded Systems}
}

\RequirePackage{ifpdf} \ifpdf
  \pdfoutput=1
  \pdftrue
  \message{pdfLaTeX}
  \documentclass[pdftex,12pt,a4paper,twoside,ngerman,numbers=noenddot]{scrbook}
  \usepackage{float}
  \usepackage[pdftex]{thumbpdf}
  \usepackage[pdftex]{pdflscape}
  \usepackage[pdftex]{graphicx}
  \usepackage[pdftex]{hyperref}
  \usepackage{pdfpages}
  \pdfoutput=1
  \pdfcompresslevel=9
  \DeclareGraphicsExtensions{.pdf,.jpg,.png}
\else
  \pdffalse
  \message{LaTeX}
  \documentclass[dvips,12pt,a4paper,twoside,ngerman,numbers=noenddot]{scrbook}
  \usepackage{float}
  \usepackage{graphicx}
  \usepackage{epsf}
  \usepackage[dvips]{hyperref}
  \DeclareGraphicsExtensions{.eps}
\fi


%
\hypersetup
{
    pdfauthor = {\Author},
    pdftitle = {\Thesistitle},
    pdfsubject = {\Thesistyp, TU Dortmund, Fakultät Informatik},
    pdfproducer = {LaTeX},
    pdfview = FitV,
    pdfstartview = FitV,
    pdfhighlight = /I,
    pdfborder = 0 0 0,
    colorlinks = false,
    bookmarksopen,
    bookmarksopenlevel = 1,
    bookmarksnumbered = false,
    plainpages = false
}%


%
\usepackage[a4paper,left=3.5cm,right=2.5cm,bottom=3.5cm,top=3cm]{geometry}
\setlength{\headheight}{15pt}
% -------------------------------------------------------------------
%
\usepackage{amsmath,amssymb}
%\usepackage{flafter}
\usepackage{subfigure}

% -------------------------------------------------------------------
\usepackage{ifthen}
\usepackage[colorinlistoftodos,prependcaption]{todonotes}

% -------------------------------------------------------------------
\usepackage[absolute,overlay]{textpos}
\setlength{\TPHorizModule}{1mm}
\setlength{\TPVertModule}{\TPHorizModule}
\textblockorigin{0mm}{0mm}
\usepackage{fix-cm}
\usepackage{setspace}
\usepackage{scrhack}
% -------------------------------------------------------------------
%
\usepackage[german,ngerman]{babel}
\usepackage[utf8]{inputenc}
\usepackage[T1]{fontenc}
\usepackage{ae,aecompl}

% -------------------------------------------------------------------
\usepackage[numbers,sort,square]{natbib}

% -------------------------------------------------------------------
\usepackage[babel,german=quotes]{csquotes}

% -------------------------------------------------------------------
\usepackage{url}
\usepackage[hyphenbreaks]{breakurl}
\def\UrlBreaks{\do\a\do\b\do\c\do\d\do\e\do\f\do\g\do\h\do\i\do\j\do\k\do\l%
\do\m\do\n\do\o\do\p\do\q\do\r\do\s\do\t\do\u\do\v\do\w\do\x\do\y\do\z\do\0%
\do\1\do\2\do\3\do\4\do\5\do\6\do\7\do\8\do\9\do\-}%

% -------------------------------------------------------------------
\usepackage[margin=0pt,font=small,labelfont=bf]{caption}

% -------------------------------------------------------------------
\usepackage{booktabs}

% -------------------------------------------------------------------
\usepackage{eurosym}

% -------------------------------------------------------------------
\renewcommand{\baselinestretch}{1.25}
\renewcommand{\topfraction}{0.9}
\renewcommand{\bottomfraction}{0.8}

% -------------------------------------------------------------------
%\clubpenalty = 10000
%\widowpenalty = 10000 \displaywidowpenalty = 10000

\parindent=0cm


% -------------------------------------------------------------------
\usepackage{fancyhdr}
\usepackage{extramarks}

\pagestyle{fancy}
\renewcommand{\chaptermark}[1]{\markboth{#1}{}}
\renewcommand{\sectionmark}[1]{\markright{#1}{}}

\fancyhf{}
\fancyhead[LE,RO]{\thepage}
\fancyhead[RE]{\textit{\nouppercase{\leftmark}}}
\fancyhead[LO]{\textit{\nouppercase{\rightmark}}}

\fancypagestyle{plain}{ %
\fancyhf{} % remove everything
\renewcommand{\headrulewidth}{0pt} % remove lines as well
\renewcommand{\footrulewidth}{0pt}} \pagestyle{headings}



% -------------------------------------------------------------------
\usepackage{color}
\definecolor{TUGreen}{rgb}{0.517,0.721,0.094}
\definecolor{TUOrange}{rgb}{1.0,0.7176,0.0}
\definecolor{BrightGray}{gray}{0.9}
\definecolor{DarkGray}{gray}{0.2}
\definecolor{white}{rgb}{1,1,1}
\definecolor{black}{rgb}{0,0,0}
\definecolor{red}{rgb}{1,0,0}




% -------------------------------------------------------------------
\usepackage{listings}

\lstdefinestyle{C++}
{
language=C++,
backgroundcolor=\color{BrightGray},
keywordstyle=\texttt\bfseries,  %\color{TUGreen}\bfseries,
commentstyle=\color{DarkGray},
stringstyle=\color{red},
showstringspaces=false,
basicstyle=\small\color{black},
numbers=left,
captionpos=b,
tabsize=4,
breaklines=true
}


% -------------------------------------------------------------------
% Algorithmen
\usepackage[plain,chapter]{algorithm}
\usepackage{algorithmic}

\usepackage{enumerate}

% -------------------------------------------------------------------
% Algorithmen anpassen
\renewcommand{\algorithmicrequire}{\textit{Eingabe:}}
\renewcommand{\algorithmicensure}{\textit{Ausgabe:}}
\floatname{algorithm}{Algorithmus}
\renewcommand{\listalgorithmname}{Algorithmenverzeichnis}
\renewcommand{\algorithmiccomment}[1]{\color{grau}{// #1}}


% -------------------------------------------------------------------
% -------------------------------------------------------------------
% -------------------------------------------------------------------
\begin{document}


% Front Page ---------------------------------------------------------
%
\pdfbookmark{Front page}{pdf:title}
\pagenumbering{alpha}
\pagestyle{empty}
%========================================================================================
% TU Dortmund, Computer Science VII
%========================================================================================

\begin{titlepage}

\begin{textblock}{150}(30.5,10.75)%
\raggedright
\includegraphics[width=83.25mm]{images/tud_logo_cmyk.pdf}%
\end{textblock}

\begin{textblock}{150}(21.2,41.6)%
\raggedright\sffamily%\Huge
{\color{red}\rule{5mm}{5mm}}
\end{textblock}

\begin{textblock}{150}(30.4,38)%
\raggedright
\includegraphics[height=13mm]{images/fi_text.pdf}
\end{textblock}

\begin{textblock}{89}(35.0,62.75)%
\begin{minipage}{80mm}
	\vfill
	\begin{center}
	\fontsize{24pt}{24pt} \sffamily
	\Thesistyp
	
	\vspace{1cm}
	\begin{onehalfspace}
    \fontsize{18pt}{18pt}
    \sffamily \Thesistitle
    \end{onehalfspace}
	
	\vspace{12mm}
\begin{onehalfspace}
	{\fontsize{14pt}{14pt}\sffamily \Author

	\today}
 \end{onehalfspace}
	\end{center}
	\vfill
\end{minipage}\end{textblock}

\begin{textblock}{150}(44.25,208)%
\begin{minipage}{120mm}
	\large
	\raggedright
	\sffamily
    {\fontsize{14pt}{14pt}
	\textbf{Supervisors:}\\
	\FirstSupervisor\\
	\SecondSupervisor\\}
\end{minipage}
\end{textblock}



\begin{textblock}{150}(44.25,242.0)%
\ifthenelse{\boolean{boolNoSecondChair}}{
%
\begin{minipage}{120mm}
	\fontsize{11.75pt}{11.75pt}\selectfont
	\raggedright
	\sffamily
	\textcolor{TUGreen}{\FirstChair}\\
	\textcolor{TUGreen}{\FirstChairTitle}\\
    \textcolor{TUGreen}{TU Dortmund}
\end{minipage}
%
}{
%
\begin{tabular*}{\textwidth}[t]{c c}%
  \begin{minipage}[t]{70mm}
    \raggedright
	\sffamily
    \textcolor{TUGreen}{\SecondChair}\\
    \textcolor{TUGreen}{(\SecondChairTitle)}\\
    \textcolor{TUGreen}{TU Dortmund}
    \end{minipage}
    \hspace*{0.5cm}
    \begin{minipage}[t]{70mm}
    \raggedright
	\sffamily
    \textcolor{TUOrange}{\SecondChair}\\
    \textcolor{TUOrange}{(\SecondChairTitle)}\\
    \textcolor{TUOrange}{TU Dortmund}
  \end{minipage}
\end{tabular*}
%
}
%
\end{textblock}



\vspace*{20cm}



\end{titlepage}


\pagestyle{empty} \cleardoublepage

% Inhaltsverzeichnis -------------------------------------------------
%
\pdfbookmark{Inhaltsverzeichnis}{pdf:toc}
\tableofcontents

\cleardoublepage \pagestyle{headings}

% Mathematische Notation -----------------------------------------------
%
\pagestyle{empty}
\pdfbookmark{Mathematische Notation}{pdf:Notation}
%========================================================================================
% TU Dortmund, Computer Science VII
%========================================================================================

\chapter*{Mathematische Notation} \label{Notation}

\newcommand{\tabdummy}{\midrule[0pt]}

\begin{tabular}{p{0.25\textwidth}p{0.65\textwidth}}
  \textbf{Notation} & \textbf{Bedeutung} \\ \toprule[1pt]
   $\mathbf{x_{i}}$ & Eingabe in das Netz \\ \tabdummy
   $\mathfrak{B}$ & Menge der Batches des Datensatzes $\mathfrak{B}=\left\{ \mathcal{B}_{1} , \ldots \mathcal{B}_{k}   \right\} $\\ \tabdummy
   $\mathcal{B}$ & ein Batch mit $m$ Elemente $\mathcal{B}=\left\{ \mathbf{x}_1, \ldots, \mathbf{x}_m \right\} $\\ \tabdummy
   $\mathbf{u}_{i,j}$ & Ausgabe der Faltungsschicht j wobei in das Netz $\mathbf{x}_i$ eingegeben wurde.\\ \tabdummy
   $\mathbf{v}_{i,j}$ & Ausgabe der Batchnormalisierungsschicht j wobei in das Netz $\mathbf{x}_i$ eingegeben wurde.\\ \tabdummy
   $\mathcal{W}$ & Gewichte der Schichten $\mathcal{W}= \left\{ W_1, \ldots, W_J \right\} $ geordnet nach Stelle der Schicht im Vorwärtsdurchgang. \\ \tabdummy
   $\mathcal{W}^k$ & Gewichte des k-ten Trainingsdurchlaufs  \\ \tabdummy
   $f\left( \mathbf{x_i}, \mathcal{W}\right) $ & Funktion, die das CNN berechnet. $f\left( \mathbf{x}_i, \mathcal{W}\right)$ berechnet eine Klassenzugehörigkeit für $\mathbf{x}_i$ \\ \tabdummy
   $y_i$ & tatsächliche Klasse von $\mathbf{x}_i$ \\ \tabdummy     
   $l\left(f( \mathbf{x_i}, \mathcal{W}\right), y_i) $ & Verlustfunktion, die durch das CNN minimiert wird
\end{tabular}

\cleardoublepage

\pagenumbering{arabic}
\pagestyle{fancy}

% Inhalte --------------------------------------------------------------
%
\chapter{Einleitung}
\label{sec:EinleitungGesamt}

\section{Motivation und Hintergrund dieser Arbeit}


MorphNet ist eine Möglichkeit mittels Verkleinern und Vergrössern des Netzwerkes effizient die Struktur eines Netzwerks zu lernen. Ist dieses Konzept auch auf PruneTrain zu übertragen und so zu erweitern, dass nicht das komplette Netz erweitert wird, sondern eben nur sinnvolle Bereiche?
Ist dieser Prozess mit Hilfe weiterer Methoden beschleunigbar?

\todo[inline]{Einleitung fertigschreiben -- zum Schluss}

\section{Aufbau der Arbeit}
\todo[inline]{Aufbau der Arbeit -- erst bei fortgeschrittener Arbeit schreiben}


\section{Suchbegriffe}
\todo[inline]{verwendete Suchbegriffe}


\cleardoublepage
\chapter{Einleitung}
\label{sec:EinleitungGesamt}

\section{Motivation und Hintergrund dieser Arbeit}


MorphNet ist eine Möglichkeit mittels Verkleinern und Vergrössern des Netzwerkes effizient die Struktur eines Netzwerks zu lernen. Ist dieses Konzept auch auf PruneTrain zu übertragen und so zu erweitern, dass nicht das komplette Netz erweitert wird, sondern eben nur sinnvolle Bereiche?
Ist dieser Prozess mit Hilfe weiterer Methoden beschleunigbar?

\todo[inline]{Einleitung fertigschreiben -- zum Schluss}

\section{Aufbau der Arbeit}
\todo[inline]{Aufbau der Arbeit -- erst bei fortgeschrittener Arbeit schreiben}


\section{Suchbegriffe}
\todo[inline]{verwendete Suchbegriffe}


\section{Verringerung der für Berechnungen nötige Zeit}

Die Zeit, die ein Convolutional Layer braucht um berechnet zu werden hängt ab von:
\todo{Fehlt hier noch etwas?}
\begin{itemize}
 \item der Filtergr\"osse
 \item der Bildgr\"osse
 \item dem verwendeten Zahlenformat
\end{itemize}
Beim Verändern der Filter- oder der Bildgr\"osse, um Trainingszeit zu sparen, ver\"andert sich auch die Erkennungsleistung \todo{cite}. Dies ist beim Verändern des verwendeten Zahlenformats nicht umbedingt gegeben. Standardformat ist eine 32 Bit Gleitkommazahl. Die einfachste Methode hier Trainingszeit zu sparen ist das Halbieren der Bitanzahl auf 16 Bit. Eine weitere Methode ist das Benutzen von 16 Bit Dynamischen Festkommazahlen.
Die beiden alternativen Methoden haben unterschiedliche Anforderungen an die Ausführungsplattform. Diese Anforderungen und die Besonderheiten der beiden Verfahren werden in den folgenden zwei Unterkapiteln näher beleuchtet.


\subsection{Berechnung mit 16 Bit Gleitkomma}

Die 16 Bit Gleitkommazahl unterscheidet sich nicht nur in der Länge von der 32 Bit Zahl sondern aus der unterschiedlichen Länge erwachsen Unterschiede in den darstellbaren Zahlen. In Tabelle \todo{ref} sind diese Unterschiede dargestellt.





Diese Nachteile von 16 Bit Gleitkommazahlen können durch drei Techniken abgemeildert oder sogar komplett aufgehoben werden:
\begin{itemize}
 \item 32 Bit Mastergewichte und Updates
 \item Sklaierung der Loss-Funktion
 \item Arithmetische Präzision 
\end{itemize}

Diese drei Techniken werden in den drei folgenden Unterkpaiteln behandelt.

\subsubsection{32 Bit Mastergewichte und Updates}

Beim Trainieren von neuronalen Netzwerken mit 16 Bit Gleitkommazahlen werden die Gewichte, Aktivierungen und Gradienten im 16 Bit Format gespeichert. Die Speicherung der Gewichte als 32 Bit Mastergewichte hat zwei mögliche Erklärungen, die aber nicht immer zutreffen müssen. 

Um nach einem Forward Druchlauf des Netzes die Gewichte abzudaten wird ein Gradientenabstiegsverfahren benutzt. Hierbei werden die Gradienten der Gewichte berechnet. Um für die Funktion, die das CNN approximiert einen besseren Approximationserfolg zu erlangen wird dann dieser Gradient mit der Lernrate multipliziert. Wird dieses Produkt in 16 Bit abgespeichert, so ist in viele Fällen das Produkt der beiden Zahlen gleich Null. Dies liegt an der Taqtsache, dass wie in Tabelle \todo{ref} zu sehen ist die kleinste darstellbare Zahl in 16 Bit wesentlich grösser ist als in 32 Bit.


Der zweite Grund wieso man Mastergewichte brauchen könnte ist die Tatsache, dass bei grossen Gewichten die Länge der Mantisse nicht ausreicht, um sowohl das Gewicht als auch das zu  addierende Update zu speichern.

Aus den beiden Gründen wird das in Abbildung \todo{ref} gezeigte Schema zum Trainieren einer Schicht mit gemischt präzisen Gleitkommazahlen benutzt.

\missingfigure{Schema}

\subsubsection{Sklaierung der Loss-Funktion}

\subsubsection{Arithmetische Präzision}


\subsection{Berechnung mit 16 Bit Dynamischen Festkommazahlen}


Quelle: \cite{FPGpu}


\subsection{Beschleunigung der Berechnung des Gradientenabstiegsverfahren}
\todo[inline]{Überblick schreiben; 4 Stunden}
Bei der Beschleunigung der Berechnung des Gradientenabstiegsverfahren gibt es vier verschiedene publizierte Herangehensweisen:
\begin{itemize}
 \item Accelerating CNN Training by Sparsifying Activation Gradients
 \item Weight Normalization: A Simple Reparameterization
to Accelerate Training of Deep Neural Networks
 \item Accelerating Deep Neural Network Training with Inconsistent Stochastic Gradient Descent
 \item Accelerated CNN Training Through Gradient Approximation 
\end{itemize}


\subsubsection{Accelerating CNN Training by Sparsifying Activation Gradients}

\subsubsection{Weight Normalization: A Simple Reparameterization
to Accelerate Training of Deep Neural Networks}




\subsubsection{Accelerating Deep Neural Network Training with Inconsistent Stochastic Gradient Descent}



\subsubsection{Accelerated CNN Training Through Gradient Approximation }




\section{Verfahren um weniger Trainingsdaten zu verwenden}

\subsection{Stochastisches Pooling}




\subsection{Lernen von Struktur und Stärke von CNNs}


\chapter{Strukturelle Veränderungen zur Beschleunigung des Trainings}

\section{Pruning um Trainingszeit zu minimieren}
Pruning ist eine Technik, die entwickelt wurde, um die Inferenzzeit eines neuronalen Netzwerks zu reduzieren. Das Pruningverfahren wird auf das bereits trainierte Netz angewendet. Dabei wird entschieden, welche Gewichte nur einen minimalen Effekt auf das Klassifikationsergebnis haben um diese zu entfernen.

Aktueller Gegenstand der Forschung ist hier die Frage, ob diese kleineren Netzwerke nicht bereits ab Epoche Null trainiert werden können, um so Trainingszeit zu sparen. Dieser Ansatz wurde in verschiedenen Veröffentlichungen untersucht:
\begin{itemize}
 \item Prune Train 
 \item The Lottery Ticket Hypothesis
\end{itemize}
Zunächst werden die einzelen Verfahren erläutert, um sie danach miteinander zu vergleichen.

\subsection{Prune Train}
Prune Train fügt einen Normalisierungsterm zur Loss-Funktion des Netzwerkes hinzu. Dies geschieht, damit der Optimierungsprozess dazu gezwungen wird möglichst kleine Gewichte zu wählen. Durch diesen Prozess wird aus dem dense Netz ein sparse Netz. Dieses sparse Netz sorgt allerdings noch nicht für weniger Zeitbedarf einer Trainingsepoche, da für ein Sparse aufwändige Datenindextechniken notwendig sind. Daher wird bei diesem Verfahren das Netz rekonfiguiert um das Modell kleiner und die Struktur wieder dense zu machen.
Dabei hat Prune Train drei zentrale Optimierungsverfahren:
\begin{itemize}
   \item eine systematische Methode zur Berechnung des group lasso Regularisierung Sanktions Koeffizienten beim Beginn des Trainings.
   \item Kanal union, ein Speicheraufruf kosteneffizientes und Index-freies Kanal Pruning Verfahren für moderne CNNs mit Kurzschlussverbindungen.
   \item Ein dynamische Mini-Batch Adjustment, dass die Größe des Mini-Batch anpasst. Dies geschieht durch beobachten des Speicherkapazitätgebrauchs einer Trainingsiteration nach jeder Pruning reconfiguration.
\end{itemize} 

Der group lasso Regularisierung Sanktions Koeffizienten ist ein Hyperparameter, der einen Trade-off  zwischen der Modellgröße und der Accuracy bildet.
Voherige Arbeiten suchen nach einem geeignetem Sanktionsmaß, was das Einbeziehen des Prunings vom Anfang des Trainings sehr teuer macht.
Unser Mechanismus kontrolliert die Group lasso Regularisierungstärke und erreicht eine hohe Modellpruningrate mit nur einem kleinen Einfluss auf die Accuracy bei nur einem Trainingsdurchlauf.
Kurzschlussverbindungen werden in modernen CNNs häufig genutzt.
Prunning aller genullten Kanäle solcher CNNs brauchen regelmässige Tensor Umordnung um die Kanalindizes zwischen den Schichten zu matchen. Dies vermindert die Performance.
Diese Umordnung wird durch den Channel Union Algorithmus vermieden. Daher folgt eine 1.9 fache Beschleunigung des Convolutional Layetrs.

Dynmaisches Mini Batch Adjustment kompensiert die verminderte Datenparallelität aufgrund des kleineren geprunten Modells durch Erhöhung der Mini-Batch Größe.
Dies sorgt sowohl für bessere Ausnutzung der Hardware ressourcen als auch zur Reduzierung des KOmmunikation overheads durch eine Verminderte Modell Update Frequenz. Beim Erhöhen der Mini-Batch Größe wird auch die Lernrate mit demselben Verhältnis erhöht, um die Accuracy nicht zu verändern.
 $$ \underset{min}{W} \left( \frac{1}{N} \sum_{i=1}^{N} l(y_i,f(x_i, W)) + \sum_{g=1}^{G} \lambda_g \cdot || W_g ||_2 \right) $$
 \begin{itemize}
  \item $ \frac{1}{N} \sum_{i=1}^{N} l(y_i,f(x_i, W))$ Standard-Kreuzentropie
  \item $\sum_{g=1}^{G} \lambda_g \cdot || W_g ||_2 $ group lasso Regulierungsterm
  \item $f(x_i)$ Vorhersage des Netzwerks auf Eingabe $x_i$
  \item $W$ Gewichte
  \item $l$ Verlustfunktion der Klassifikation und Grundwahrheit $y_i$
  \item $N$ Minibatchgröße
  \item $G$ Zahl von Gruppen
  \item $\lambda$ Verdünnungskoeffizient
\end{itemize}

$$ \lambda \cdot \sum_{l=1}^{L} \left( \sum_{c_l=1}^{C_l} || W_{c_l,:,:,:} ||_2 + \sum_{k_l=1}^{K_l} || W_{:,k_l,:,:}||_2 \right) $$ 
Design eines speziellen Groupo Lasso Regulierers, der Gewichte jedes Kanals (Input oder Output) und jeder Schicht gruppiert. $\lambda$ wird als einziger globaler Regularisierungsfaktor gewählt, da so der Fokus auf dem Vermindern der Rechenzeit liegt und nicht auf der Modellgröße. Dies hat zur Folge, dass vorallem große Features verdünnt werden, was zu einer größeren Verminderung der Rechenleistung führt.

Um die Lasso Group Regularisierung vom Anfang des Trainings zu benutzen sollteder Koeffizient $\lambda$ sinnvoll gewählt werden. Dies sorgt für eine hohe Vorhersageaccuracy und einer hohen Pruning Rate. Um zeitintensives Hyperparametertuning zu vermeiden wird hier eine neue Methode eingeführt:
 
 $$LPR=\frac{\lambda \sum_{g}^{G}|| W_{g,:} ||}{l(y_i,f(x_i,W))+\lambda \sum_{g}^{G}||W_{g,:} ||)} $$
 
 Berechnet wird dies durch setzen von zufälligen Werten, mit denen die Gewichte initialisiert werden. LPR wird einmal berechnet und dann bis zum Ende weiter benutzt.

 
 Nach jedem solchen Intervall werden Input und Outputkanäle die 0 sind gepruned. Um ein Missverhältnis zwischen den Dimensionen zu verhindern wird nur die Verbindung von 2 verdünnten Kanälen von 2 aufeinanderfolgenden Schichten gepruned. Alle Trainingsvariablen bleiben gleich.
 Das Reconfigurationsintervall ist ist der einzige zusätzliche Hyperparameter. Zu gross gewählt würde der Intervallparameter zu wenig Zeitverbesserung bringen. Zu klein gewählt könnte er die Lernqualität beeinflussen.
4 Matriken zur Evaluierung: Training und Inference FLOPs, gemessenen Trainingszeit, und Validierungsaccuracy. 




\section{Net 2 Net}


\section{Kernel rescaling}


\section{Resource Aware Layer Replacement}


\section{Additive Methoden}
Die Methoden in diesem Kapitel beeinflussen die Trainingszeit nicht direkt, sondern helfen die Folgen anderer Verfahren abzumildern.

\subsection{Ghost Batch Normalization}





\cleardoublepage
\chapter{Einleitung}
\label{sec:EinleitungGesamt}

\section{Motivation und Hintergrund dieser Arbeit}


MorphNet ist eine Möglichkeit mittels Verkleinern und Vergrössern des Netzwerkes effizient die Struktur eines Netzwerks zu lernen. Ist dieses Konzept auch auf PruneTrain zu übertragen und so zu erweitern, dass nicht das komplette Netz erweitert wird, sondern eben nur sinnvolle Bereiche?
Ist dieser Prozess mit Hilfe weiterer Methoden beschleunigbar?

\todo[inline]{Einleitung fertigschreiben -- zum Schluss}

\section{Aufbau der Arbeit}
\todo[inline]{Aufbau der Arbeit -- erst bei fortgeschrittener Arbeit schreiben}


\section{Suchbegriffe}
\todo[inline]{verwendete Suchbegriffe}


\chapter{Durchführung}

\begin{itemize}
 \item Baseline Training ab der ersten Epoche.
 \item Prune währenddem Training
 \item Training bis zu dem Zeitpunkt wo durch ein weitere Epoche nichts besser wird
 \item Überprüfe wieviel in den letzten Epochen gepruned wurde um zu entscheiden ob das Netz weiter odertiefersein soll
\end{itemize}



% Anhang ---------------------------------------------------------------
%
\cleardoublepage
\appendix

%========================================================================================
% TU Dortmund, Computer Science VII
%========================================================================================
\chapter{d}


%
\listoffigures
\addcontentsline{toc}{chapter}{Abbildungsverzeichnis}
\cleardoublepage

%
\listofalgorithms
\addcontentsline{toc}{chapter}{Algorithmenverzeichnis}
\cleardoublepage

%
\renewcommand{\lstlistlistingname}{Quellcodeverzeichnis}
\lstlistoflistings
\addcontentsline{toc}{chapter}{Quellcodeverzeichnis}
\cleardoublepage

%
\addcontentsline{toc}{chapter}{Literaturverzeichnis}
\bibliographystyle{alpha}
\bibliography{Literature}

% ----------------------------------------------------------------------

\end{document}
